\documentclass[11pt,a4paper]{article}
\usepackage[a4paper,margin=2.4cm]{geometry}
\usepackage{graphicx}
\usepackage{float}

\title{Labwork 2}
\author{Pham Tien Nhat \\ Student ID: 2540090 \\ Advanced HPC -- USTH}
\date{\today}

\begin{document}
\maketitle

\section{Device Name}
\texttt{cuda.detect()} finds the supported GPU; \texttt{cuda.select\_device(0)}
selects it. Device is: NVIDIA GeForce RTX 3050 Laptop GPU.

\begin{figure}[H]
\centering
\includegraphics[width=0.7\textwidth]{../../screenshots/lab2/1.png}
\caption{\texttt{cuda.detect()} showing the CUDA device.}
\end{figure}

\begin{figure}[H]
\centering
\includegraphics[width=0.55\textwidth]{../../screenshots/lab2/2.png}
\caption{\texttt{cuda.select\_device(0)} returning device id and name.}
\end{figure}

\section{Core Info}
Compute capability 8.6 (Ampere), 16 multiprocessors.

\begin{figure}[H]
\centering
\includegraphics[width=0.55\textwidth]{../../screenshots/lab2/3.png}
\caption{Compute capability and multiprocessor count from device.}
\end{figure}

\section{Memory Info}
\texttt{get\_memory\_info()} reports free and total device memory in bytes.


\begin{figure}[H]
\centering
\includegraphics[width=0.6\textwidth]{../../screenshots/lab2/4.png}
\caption{Free and total memory returned by \texttt{get\_memory\_info()}.}
\end{figure}

\end{document}
